\providecommand{\kBKM}{\kappa_{\mathrm{BKM}}}
\providecommand{\kch}{\kappa_{\mathrm{ch}}}
\providecommand{\kcat}{\kappa_{\mathrm{cat}}}
\providecommand{\Mbar}{\overline{M}}
\providecommand{\Nred}{\mathsf{N}}
\providecommand{\Zred}{\mathcal{Z}}
\providecommand{\PhiOP}{\Phi_{10}^{\mathrm{OP}}}
\providecommand{\Phiun}{\Phi_{10}^{\mathrm{un}}}
\providecommand{\DeltaFive}{\Delta_5}
\providecommand{\DFive}{D_5}

\section*{Oberdieck-Pixton and the primitive \texorpdfstring{$K3\times E$}{K3 x E} scalar}

\subsection*{Primitive reduced Gromov-Witten theorem}

Let $S$ be a non-singular projective K3 surface, let $E$ be an elliptic
curve, and set $X=S\times E$.  For a primitive class
$\beta_h\in H_2(S,\mathbb Z)$ with
$\langle\beta_h,\beta_h\rangle=2h-2$, and for $d\geq0$, Oberdieck and
Pixton define the reduced invariant
\[
  \Nred_{g,\beta_h,d}^{X}
  =
  \int_{{\left[\Mbar_{g,1}^{\bullet}(X,(\beta_h,d))\right]}^{\mathrm{red}}}
  \operatorname{ev}_1^*\!\left(\pi_1^*(\beta_h^\vee)\cup
  \pi_2^*([\mathrm{pt}])\right).
\]
The reduced virtual class has dimension $1$.  The insertion kills the
one-dimensional $E$-translation orbit.  The invariant is therefore a
reduced Gromov-Witten count modulo translation, not ordinary
cohomology, reduced homology, or a motivic reduction.

Deformation invariance makes the invariant depend only on $g,h,d$.
Set
\[
  \Zred(u,q,\tilde q)
  =
  \sum_{g\geq0}\sum_{h\geq0}\sum_{d\geq0}
  \Nred_{g,h,d}^{X}\,u^{2g-2}q^{h-1}\tilde q^{d-1}.
\]
Oberdieck-Pixton, \emph{Holomorphic anomaly equations and the Igusa
cusp form conjecture}, Invent.\ Math.\ \textbf{213} (2018), no.~2,
pages 507 to 587, prove
\[
  \boxed{\;\Zred(u,q,\tilde q)
  =
  -\frac{1}{\chi_{10}(p,q,\tilde q)}\;,\qquad p=e^{iu}.}
\]
The equality is the Laurent expansion in the chamber
$|p^{k}q^{h}\tilde q^{d}|<1$ whenever $4hd-k^2\geq -1$.

The product convention is the monic Igusa convention
\[
  \chi_{10}(p,q,\tilde q)
  =
  pq\tilde q
  \prod_{k,h,d}
  {(1-p^{k}q^{h}\tilde q^{d})}^{c(4hd-k^2)},
\]
where
\[
  \sum_n c(n)q^n
  =
  \frac{40{\theta_3(q)}^4-8{\theta_4(q)}^4}
       {\theta_3(q){\theta_2(q)}^4}
  =
  2q^{-1}+20-128q^3+216q^4-\cdots .
\]
The product ranges over all $k\in\mathbb Z$ and $h,d\geq0$ satisfying
one of the three inequalities
\[
  h>0,\qquad
  h=0,\ d>0,\qquad
  h=d=0,\ k<0.
\]

\subsection*{Normalisation by \texorpdfstring{$\Delta_5$}{Delta5}}

The primitive Borcherds denominator is the weight-$5$ Gritsenko-Nikulin
form $\DeltaFive$.  With the monic product convention, its first
Fourier coefficient is divided by $64$:
\[
  \DFive=64^{-1}\DeltaFive .
\]
Hence the Oberdieck-Pixton scalar form is
\[
  \PhiOP
  =
  \DFive^2
  =
  4096^{-1}\DeltaFive^2,
  \qquad
  \Phiun=\DeltaFive^2.
\]
Consequently the scalar partition function is
\[
  Z^{\mathrm{red}}_{GW,\mathrm{prim}}(K3\times E)
  =
  -\PhiOP^{-1}
  =
  -\DFive^{-2}
  =
  -4096\,\DeltaFive^{-2}.
\]
The coefficient $4096=64^2$ is the square of the theta-normalised to
monic-product conversion.  It records neither an anomaly nor an
orientation, degeneracy, torsion, or CHL factor.

The primitive BKM weight belongs to the square root:
\[
  \kBKM(\mathfrak g_{\DeltaFive})
  =
  \kBKM(\DeltaFive)
  =
  c_1(0)/2
  =
  10/2
  =
  5.
\]
The Oberdieck-Pixton theorem proves the scalar reciprocal square.  The
BKM Lie superalgebra, Pfaffian square root, orientation character, and
Hall-Borcherds recognition map are additional structures.

\subsection*{Stable-pairs reading}

The displayed theorem is a reduced Gromov-Witten statement.  Its
stable-pairs reading is legitimate in the primitive K3 class because
the proof uses the reduced GW/PT correspondence: the reduced
Pandharipande-Thomas invariant is computed by integrating the Behrend
function on $P_n(X,(\beta_h,d))/E$, equivalently by the reduced virtual
class.  Thus the precise DT/PT notation is
\[
  Z^{\mathrm{red}}_{PT,\mathrm{prim}}(K3\times E)
  =
  -\PhiOP^{-1}.
\]
An unreduced Donaldson-Thomas invariant, an ordinary virtual class, and
a quotient lacking the $E$-translation reduction belong to different
theorems.

\subsection*{Boundary of the theorem}

The theorem is an $N=1$ primitive reduced result.  The class
$\beta_h$ is primitive in $H_2(S,\mathbb Z)$.  Classes
$\beta=m\beta_0$ for $m>1$ require separate theorems.
Oberdieck-Pandharipande's all-class Conjecture~B expresses imprimitive
series by divisor sums over primitive series.
That divisor sum is a proof obligation beyond the Oberdieck-Pixton
theorem.

Finite quotients, torsion sectors, and CHL extensions require:
\[
  \text{quotient orientation line},\qquad
  \text{descent of the reduced obstruction theory},\qquad
  \text{compatibility with the }E\text{-translation quotient},
\]
and the stabiliser-weighted divisor sums in the chosen chamber.  These
data are exactly the proof of a replacement
\[
  -\PhiOP^{-1}
  \leadsto
  \text{CHL/torsion/finite-quotient scalar}
\].

\subsection*{Citation and formula normalisations}

Every citation to this result must name both Oberdieck and Pixton.  The
bibliographic record is:
\[
  \text{Invent.\ Math.\ \textbf{213} (2018), no.~2, pages 507 to 587,}
  \qquad
  \mathrm{doi}:10.1007/s00222-018-0794-0.
\]
The common citation ``Invent.\ Math.\ \textbf{222} (2020)'' belongs to a
different bibliographic record.

The normalisation chain is:
\[
  \boxed{\;
  \chi_{10}^{\mathrm{OP}}=\PhiOP=\DFive^2
  =4096^{-1}\DeltaFive^2,\qquad
  Z^{\mathrm{red}}_{GW,\mathrm{prim}}
  =Z^{\mathrm{red}}_{PT,\mathrm{prim}}
  =-\PhiOP^{-1}
  =-4096\,\DeltaFive^{-2}.
  \;}
\]
The unnormalised equality $\Phiun=\DeltaFive^2$ is harmless only when it
is explicitly separated from the Oberdieck-Pixton scalar convention.
The Oberdieck-Pixton reduced scalar is
\(-4096\,\DeltaFive^{-2}\).

\subsection*{References}

Oberdieck-Pixton 2018: \emph{Holomorphic anomaly equations and the
Igusa cusp form conjecture}, arXiv:1706.10100, Invent.\ Math.\
\textbf{213} (2018), no.~2, pages 507 to 587.

Oberdieck-Pandharipande 2016: \emph{Curve counting on $K3\times E$, the
Igusa cusp form $\chi_{10}$, and descendent integration},
arXiv:1411.1514, in \emph{K3 Surfaces and Their Moduli}, Progress in
Mathematics \textbf{315}, pages 245 to 278.

Gritsenko-Nikulin 1998: \emph{Automorphic forms and Lorentzian
Kac-Moody algebras. Part II}, Internat.\ J.\ Math.\ \textbf{9},
pages 201 to 275.
